% Figure: the three damping regimes of a second-order RLC step response
% located on the s-plane by the characteristic roots, with the matching
% time-domain step responses. Left: root loci; right: step curves. pgfplots.
\begin{figure}[htbp]
\centering
% === left: s-plane root positions ===
\begin{tikzpicture}
\begin{axis}[
  width=6.2cm, height=6.2cm,
  xlabel={$\Re(s)$ ($\times 10^{4}\,\si{\per\second}$)},
  ylabel={$\Im(s)$ ($\times 10^{4}\,\si{\per\second}$)},
  xmin=-12, xmax=2, ymin=-7, ymax=7,
  axis lines=middle,
  tick label style={font=\tiny},
  label style={font=\tiny},
  clip=false,
]
% underdamped pair (complex)
\addplot[only marks, mark=x, engblue, mark size=3pt]
  coordinates {(-1,5.916) (-1,-5.916)};
\node[engblue, font=\tiny, anchor=south west] at (axis cs:-1,5.916) {underdamped};
% critically damped (repeated real)
\addplot[only marks, mark=square*, engorange, mark size=2pt]
  coordinates {(-6,0)};
\node[engorange, font=\tiny, anchor=south] at (axis cs:-6,0.5) {critical};
% overdamped (two distinct real)
\addplot[only marks, mark=*, engred, mark size=2pt]
  coordinates {(-1.0,0) (-11.0,0)};
\node[engred, font=\tiny, anchor=north] at (axis cs:-8,-0.6) {overdamped};
\end{axis}
\end{tikzpicture}
\hspace{0.6cm}
% === right: step responses ===
\begin{tikzpicture}
\begin{axis}[
  width=6.2cm, height=6.2cm,
  xlabel={$\omega_n t$}, ylabel={$v_C/V_s$},
  xmin=0, xmax=12, ymin=0, ymax=1.6,
  axis lines=left,
  tick label style={font=\tiny},
  label style={font=\tiny},
  legend style={font=\tiny, at={(0.98,0.03)}, anchor=south east},
  clip=false,
]
% underdamped zeta=0.17: omega_d = sqrt(1-zeta^2)=0.9854, phase=acos(zeta)=80.21 deg
\addplot[engblue, very thick, domain=0:12, samples=200]
  {1 - exp(-0.17*x)/0.9854*sin(deg(0.9854*x) + 80.21)};
\addlegendentry{$\zeta=0.17$}
% critically damped zeta=1
\addplot[engorange, thick, domain=0:12, samples=200]
  {1 - (1 + x)*exp(-x)};
\addlegendentry{$\zeta=1$}
% overdamped zeta=2: roots -0.268, -3.732 (units of omega_n)
\addplot[engred, thick, domain=0:12, samples=200]
  {1 - (3.732*exp(-0.268*x) - 0.268*exp(-3.732*x))/3.464};
\addlegendentry{$\zeta=2$}
\draw[enggray, dashed] (axis cs:0,1) -- (axis cs:12,1);
\end{axis}
\end{tikzpicture}
\caption[Second-order RLC roots and step response]{The characteristic
roots of a series RLC circuit (left, on the $s$-plane) and the step
response each root configuration produces (right). A complex-conjugate
pair off the real axis (underdamped) gives a ringing response that
overshoots; a repeated real root (critically damped) gives the fastest
monotone approach; two distinct real roots (overdamped) give a slower
sum of two decaying exponentials. The real part of the dominant root
sets the envelope decay rate; the imaginary part sets the ringing
frequency $\omega_d$.}
\label{fig:vol04ch09:rlc-roots}
\end{figure}
